Let (\mu) denote Lebesgue measure on (\mathbb R^2). For three points (x_1,x_2,x_3\in\mathbb R^2), define the (unsigned) Euclidean area of the triangle with these vertices by
[
\operatorname{Area}(x_1,x_2,x_3);=;\frac12,\bigl|\det(x_2-x_1,;x_3-x_1)\bigr|,
]
where (\det((a,b),(c,d)):=ad-bc).

The Erdős problem asks:

> Does there exist a constant (c>0) such that for **every** Lebesgue measurable set (A\subseteq\mathbb R^2) with (\mu(A)\ge c), there exist (x_1,x_2,x_3\in A) with (\operatorname{Area}(x_1,x_2,x_3)=1)?

Stress points / edge cases:

* If (\mu(A)=\infty) the statement is known to be true (per Erdős; see context below).
* The nontrivial regime is (0<\mu(A)<\infty), and (as the examples show) (A) must be allowed to have empty interior and be highly disconnected.
* A necessary lower bound for any such (c) is
  [
  c_0=\frac{4\pi}{\sqrt{27}}\approx 2.418,
  ]
  witnessed by a disk of slightly smaller area (details proved in §4).
